% !TeX root = ../main.tex
% Section 4: Duality Hedging and the Semiclassical Expansion
% Batch #4.4 — Post M1+M2 axis, notation hygiene, γ sign corrected in Remark 4.12

\section{Duality Hedging and the Semiclassical Expansion}\label{sec:duality-hedging}

Section~\ref{sec:no-arbitrage} established the \emph{existence} and uniqueness
of the optimal hedge as an orthogonal projection in the GNS Hilbert space. This
section provides its \emph{explicit form} and derives the empirically testable
consequences. Sections~4.1--4.2 show that the projection is
a quantum measurement, giving the State--Hedge pricing representation.
Section~\ref{sec:semiclassical} then expands the pricing formula in powers of
$\hbar_{\mathrm{econ}}$, producing the $\mathcal{O}(\tau^{-1})$ short-maturity
skew divergence and the quantum optimal hedge---the two theoretical predictions
that drive the empirical tests of Section~\ref{sec:empirical-calibration}.

\subsection{The Duality Between Hedging and Measurement}

The central duality of the quantum framework equates the optimal hedge for a European
option to a quantum measurement of the payoff operator in the GNS representation,
yielding a computationally tractable pricing approach.

\begin{definition}[Hedging operator]\label{def:hedging-operator}
Let $\hat{G} = g(\hat{S}_T)$ be the payoff operator for a European option with payoff
function $g$ and let $\rho \in \mathcal{T}(\mathcal{H})$ be the market state. The
\emph{optimal hedge} is the operator $\hat{\Pi} \in \mathcal{A}$ that minimizes the
quadratic hedging error
\begin{equation}\label{eq:optimal-hedge-def}
    \hat{\Pi}^* = \arg\min_{\hat{\Pi} \in \mathcal{A}}
    \operatorname{Tr}\!\left[\rho \, (\hat{G} - \hat{\Pi})^2\right],
\end{equation}
where $\mathcal{A}$ is the observable algebra generated by $\{\hat{X}, \hat{P}\}$.
\end{definition}

The optimization \eqref{eq:optimal-hedge-def} is a least-squares problem in the
non-commutative $L^2(\rho)$ space \citep{pisier2003introduction}. The solution is
the conditional expectation of $\hat{G}$ onto $\mathcal{A}$ with respect to $\rho$,
which we denote by $\mathbb{E}_\rho[\hat{G} \,|\, \mathcal{A}]$.

\begin{theorem}[Hedging--measurement duality]\label{thm:hedging-measurement}
Let $\rho \in \mathcal{T}(\mathcal{H})$ be a faithful normal state on $\mathcal{A}$
and let $\pi_\rho : \mathcal{A} \to \mathcal{B}(\mathcal{H}_\rho)$ be the GNS
representation. Then for any payoff operator $\hat{G} \in \mathcal{A}$, the optimal
hedge $\hat{\Pi}^*$ satisfies
\begin{equation}\label{eq:hedging-measurement-duality}
    \pi_\rho(\hat{\Pi}^*) |\Omega_\rho\rangle = P_{\mathcal{A}_\rho}
    \pi_\rho(\hat{G}) |\Omega_\rho\rangle,
\end{equation}
where $P_{\mathcal{A}_\rho}$ is the orthogonal projection onto the closure of
$\pi_\rho(\mathcal{A}) |\Omega_\rho\rangle$ in $\mathcal{H}_\rho$, and
$|\Omega_\rho\rangle$ is the GNS cyclic vector.
\end{theorem}

\begin{proof}
The GNS representation provides an isometric embedding of $\mathcal{A}$ into
$\mathcal{B}(\mathcal{H}_\rho)$ via $\pi_\rho$, with cyclic vector
$|\Omega_\rho\rangle$ such that
$\omega(A) = \langle \Omega_\rho \,|\, \pi_\rho(A) \Omega_\rho \rangle$ for all
$A \in \mathcal{A}$. The quadratic hedging error \eqref{eq:optimal-hedge-def} can
be written as
\begin{equation*}
    \operatorname{Tr}[\rho (\hat{G} - \hat{\Pi})^2]
    = \|\pi_\rho(\hat{G}) |\Omega_\rho\rangle
    - \pi_\rho(\hat{\Pi}) |\Omega_\rho\rangle\|_{\mathcal{H}_\rho}^2.
\end{equation*}
The minimizer is the orthogonal projection of $\pi_\rho(\hat{G}) |\Omega_\rho\rangle$
onto the closed subspace $\overline{\pi_\rho(\mathcal{A}) |\Omega_\rho\rangle}$,
which is precisely $P_{\mathcal{A}_\rho} \pi_\rho(\hat{G}) |\Omega_\rho\rangle$.
Since $\pi_\rho$ is injective (faithfulness of $\rho$), there exists a unique
$\hat{\Pi}^* \in \mathcal{A}$ satisfying \eqref{eq:hedging-measurement-duality}.
\end{proof}

\begin{remark}[Economic interpretation]\label{rem:econ-interp}
The duality theorem reveals that the optimal hedge is the \emph{quantum measurement}
of the payoff operator in the market state $\rho$. The GNS vector
$|\Omega_\rho\rangle$ represents the ``vacuum'' state of the market, and the
projection $P_{\mathcal{A}_\rho}$ extracts the component of the payoff that is
replicable by trading in the available securities. The non-replicable component,
orthogonal to $\mathcal{A}_\rho$, is the intrinsic hedging error that cannot be
eliminated by any dynamic trading strategy.
\end{remark}

\subsection{The State--Hedge Pricing Representation}

The core pricing formula is formalized as a representation theorem on a real Hilbert
space $\mathcal{H}_\omega \simeq L^2(\mathbb{R}, \mathbb{R})$ over the log-price
axis $x = \log S_T$, guaranteeing real-valued inner products and a manifestly
positive pricing functional.

\begin{definition}[Payoff-state vector]\label{def:payoff-state-vector}
Let $g(x)$ be the payoff function of a European option expressed in terms of the
log-price $x = \log S_T$, and let $\Phi_t(x)$ be the wavefunction of the market
risk state at time $t$ in the log-price representation. The \emph{payoff-state
vector} at time $t$ is the pointwise product
\begin{equation}\label{eq:payoff-state}
    Y_t(x) := g(x) \, \Phi_t(x) \in L^2(\mathbb{R}, \mathbb{R}),
\end{equation}
where the Hilbert space $\mathcal{H}_\omega \simeq L^2(\mathbb{R}, \mathbb{R})$ is
a \emph{real} $L^2$ space over the log-price axis, equipped with the standard real
inner product
\begin{equation*}
    \langle f \,|\, h \rangle = \int_{\mathbb{R}} f(x) \, h(x) \, dx.
\end{equation*}
The real nature of $\mathcal{H}_\omega$ guarantees that all duality pairings are
strictly real-valued.
\end{definition}

\begin{definition}[Hedgeable subspace at time $t$]\label{def:hedgeable-subspace-time}
Let $\mathcal{K}_t \subset \mathcal{H}_\omega$ be the closed linear subspace spanned
by the discounted asset-price wavefunctions available for hedging at time $t$:
\begin{equation*}
    \mathcal{K}_t = \overline{\operatorname{span}}
        \{\psi_{S_{t_i}}(x) : t_i \in [0, t]\},
\end{equation*}
where $\psi_{S_{t_i}}(x)$ is the wavefunction of the discounted asset price at time
$t_i$. The subspace $\mathcal{K}_t$ represents the set of all tradeable hedging
strategies.
\end{definition}

\begin{lemma}[Closedness of hedgeable subspace and market-state regularity]
\label{lem:Kt-closed}
Under the Stone--von Neumann representation $\mathcal{H}_\omega \simeq
L^2(\mathbb{R},\mathbb{R})$, $\mathcal{K}_t$ as defined above is a proper
closed subspace of $\mathcal{H}_\omega$ whenever the market depth is finite
($\hbar_{\mathrm{econ}} > 0$). Moreover, the Lindblad-evolved market risk state
satisfies $\Phi_t \in \mathcal{K}_t$ under the following regularity assumption.
\end{lemma}

\begin{assumption}[Market state in hedgeable subspace: formal statement]\label{ass:phi-in-K-detail}
The market risk state $\Phi_t$ lies in the hedgeable subspace: $\Phi_t \in
\mathcal{K}_t$ (cf.\ Assumption~\ref{ass:phi-in-K} in Section~\ref{sec:no-arbitrage}).
This holds whenever the Lindblad dynamics preserve the
asset-price sub-algebra (i.e., the market state is generated by traded
observables), as is the case for the single-qubit Bloch parametrization
of Section~\ref{sec:empirical-calibration}.
\end{assumption}

\begin{proof}[Proof of Lemma~\ref{lem:Kt-closed}]
\emph{Closedness.} By definition, $\mathcal{K}_t$ is the closed linear span
of the set $\{\psi_{S_{t_i}} : t_i \in [0,t]\}$, hence closed by construction.

\emph{Properness.} We claim $\mathcal{K}_t \subsetneq \mathcal{H}_\omega =
L^2(\mathbb{R},\mathbb{R})$ whenever the market is genuinely incomplete
($\hbar_{\mathrm{econ}} > 0$). Each wavefunction $\psi_{S_{t_i}}(x)$ is
concentrated near the asset-price level $S_{t_i}$: it belongs to $L^2(\mathbb{R})$
but its support is localized in price-space. The payoff function $g(x)$ of a
generic option has support that extends beyond the span of these wavefunctions
(e.g., for a call, $g(x) = (e^x - K)_+$ grows at $+\infty$, while any finite
linear combination of localized wavefunctions decays). Therefore $g \notin
\mathcal{K}_t$ in general, confirming $\mathcal{K}_t \subsetneq L^2(\mathbb{R})$.

\emph{Remark on incompleteness.} Note that the irreducibility of the Weyl
algebra representation (Stone--von Neumann theorem) refers to the absence of
\emph{operator-invariant} closed subspaces, not the absence of arbitrary closed
subspaces. $L^2(\mathbb{R})$ has infinitely many proper closed subspaces (e.g.,
$\{f : \mathrm{supp}(f) \subseteq [0,1]\}$), and $\mathcal{K}_t$ is one of them.
Irreducibility is not invoked here.
\end{proof}

\begin{corollary}[State--Hedge Pricing Representation, of Theorem \ref{thm:pricing}]
\label{cor:state-hedge-pricing}
Let $g(x)$ be the payoff of a European option maturing at $T$, let $r \geq 0$ be
the risk-free rate, and let $\Phi_t \in \mathcal{H}_\omega$ be the market risk
state at time $t \leq T$, evolving under the Lindblad master equation
\eqref{eq:lindblad}. Define the payoff-state vector $Y_t = g \Phi_t$ as in
\eqref{eq:payoff-state}, and let $h^* \in \mathcal{K}_t$ be the unique optimal
hedge representative given by the orthogonal projection of $Y_t$ onto the hedgeable
subspace:
\begin{equation}\label{eq:optimal-hedge-time}
    h^* := \Pi_{\mathcal{K}_t} Y_t \in \mathcal{K}_t,
\end{equation}
where $\Pi_{\mathcal{K}_t} : \mathcal{H}_\omega \to \mathcal{K}_t$ is the
orthogonal projection operator. Then by Theorem \ref{thm:pricing}, the
arbitrage-free price is given by formula \eqref{eq:pricing-master}, where
$\langle \cdot \,|\, \cdot \rangle$ is interpreted as the real inner product on
$\mathcal{H}_\omega \simeq L^2(\mathbb{R}, \mathbb{R})$, a real Hilbert space over
the log-price axis $x = \log S_T$.

\begin{remark}[Real versus complex Hilbert space]\label{rem:real-vs-complex}
The abstract GNS construction produces a \emph{complex} Hilbert space
$\mathcal{H}_\omega^{\mathbb{C}}$, as required by the spectral theory of the
Weyl operators and the Robertson inequality. The present corollary works
with its \emph{realification}
$\mathcal{H}_\omega := \mathcal{H}_\omega^{\mathbb{C}}|_{\mathbb{R}}$,
the same underlying vector space equipped with the real inner product
$\langle f,g\rangle_{\mathbb{R}} := \operatorname{Re}\langle f,g\rangle_{\mathbb{C}}$.
This is legitimate because option payoffs $g(x)$ and market risk states $\Phi_t(x)$
in the log-price representation are \emph{real-valued} functions, so the complex
inner product reduces to a real integral. The realification preserves the
Hilbert-space projection theorem (hence the existence and uniqueness of $h^*$)
and guarantees $P_t = \langle h^*|\Phi_t\rangle_{\mathbb{R}} \in \mathbb{R}_{\geq 0}$
without loss of the operator-algebraic structure of \S\ref{sec:weyl-heisenberg}.
The Stone--von Neumann theorem \citep{vonNeumann1931} ensures that the irreducible
representation of the CCR algebra on $L^2(\mathbb{R},\mathbb{C})$ is unique up
to unitary equivalence; the real subspace used here is the invariant subspace
of real-valued functions, which is preserved by the option pricing projection.
\end{remark}
\end{corollary}

\begin{proof}
By Theorem \ref{thm:pricing}, the arbitrage-free price is
$P_t = \langle h^* | \Phi_t \rangle e^{-r(T-t)}$ on the GNS Hilbert space
$\mathcal{H}_\omega$. The GNS construction (Definition \ref{def:gns}) realizes
$\mathcal{H}_\omega$ concretely as the real $L^2(\mathbb{R}, \mathbb{R})$ space
over the log-price axis $x = \log S_T$, with $\Phi_t(x)$ the wavefunction
representative of the market risk state and $Y_t(x) = g(x) \Phi_t(x)$
(Definition \ref{def:payoff-state-vector}) the payoff-state vector. The orthogonal
projection $h^* = \Pi_{\mathcal{K}_t}(e^{-rT} Y_t)$ onto the hedgeable subspace
$\mathcal{K}_t$ (Definition \ref{def:hedgeable-subspace-time}) and the
commutativity of the deterministic discount factor with the projection yield the
pricing formula \eqref{eq:pricing-master}. Realness and positivity follow
immediately from the real Hilbert space structure.
\end{proof}

\begin{remark}[Economic interpretation of the State--Hedge representation]
\label{rem:state-hedge-econ}
The State--Hedge Pricing Representation \eqref{eq:pricing-master} provides a complete
resolution of the three core economic requirements:

\paragraph{Realness.}
Because $\mathcal{H}_\omega \simeq L^2(\mathbb{R}, \mathbb{R})$ is a real Hilbert
space, the inner product $\langle h^* \,|\, \Phi_t \rangle$ is manifestly real-valued.
No complex phases appear in the pricing formula; the complex structure of the
underlying $C^*$-algebra is confined to the algebraic relations among observables
and does not affect the final price.

\paragraph{Positivity.}
The optimal hedge $h^* = \Pi_{\mathcal{K}_t} Y_t$ is the orthogonal projection of
the payoff-state vector $Y_t = g \Phi_t$ onto $\mathcal{K}_t$. Write the
decomposition
\begin{equation*}
    Y_t = h^* + \varepsilon_t, \qquad \varepsilon_t \perp \mathcal{K}_t,
\end{equation*}
where $\varepsilon_t$ is the unhedgeable residual. Since $\Phi_t \in \mathcal{K}_t$
(the market risk state is itself a tradeable quantity), the orthogonality
$\varepsilon_t \perp \Phi_t$ implies
\begin{equation*}
    \langle h^* \,|\, \Phi_t \rangle = \langle Y_t \,|\, \Phi_t \rangle
    = \int_{\mathbb{R}} g(x) |\Phi_t(x)|^2 \, dx \geq 0,
\end{equation*}
where the inequality follows from $g(x) \geq 0$ for a standard option payoff and
$|\Phi_t(x)|^2 \geq 0$. Thus the price is non-negative for any non-negative payoff,
establishing the positivity of the pricing functional.

\paragraph{Classical limit (complete markets).}
When the market is frictionless ($\hbar_{\mathrm{econ}} = 0$), the hedgeable
subspace coincides with the full space: $\mathcal{K}_t = \mathcal{H}_\omega$. In
this case $\Pi_{\mathcal{K}_t} = \mathbf{1}$, so $h^* = Y_t = g \Phi_t$, and the
pricing formula collapses to
\begin{equation*}
    P_t^{\mathrm{BSM}} = e^{-r(T-t)} \langle g \Phi_t \,|\, \Phi_t \rangle
    = e^{-r(T-t)} \int_{\mathbb{R}} g(x) |\Phi_t(x)|^2 \, dx
    = e^{-r(T-t)} \mathbb{E}^{\mathbb{Q}}[g(S_T) \,|\, \mathcal{F}_t],
\end{equation*}
where $|\Phi_t(x)|^2 = q_t(x)$ is the classical risk-neutral density. The quantum
formula therefore contains the classical BSM expectation as the special case of
complete markets.

\paragraph{Incomplete markets (general case).}
When $\hbar_{\mathrm{econ}} > 0$, the market is incomplete and $\mathcal{K}_t$ is
a proper subspace of $\mathcal{H}_\omega$. The optimal hedge $h^*$ is the best
tradeable approximation to the payoff-state vector $Y_t$, and the inner product
$\langle h^* \,|\, \Phi_t \rangle$ represents the valuation of this optimal hedge
under the market friction encoded by $\hbar_{\mathrm{econ}}$. The residual
$\varepsilon_t = Y_t - h^*$ is the intrinsic hedging error that cannot be
eliminated by any dynamic strategy---it is the cost of non-commutativity,
quantified by the economic Planck constant.
\end{remark}

\begin{corollary}[Relation to the GNS pricing formula]\label{cor:gns-pricing}
The State--Hedge representation \eqref{eq:pricing-master} is equivalent to the
abstract pricing relation \eqref{eq:master-price-intro} realized in the GNS Hilbert
space, with $|h^*\rangle = \Pi_{\mathcal{K}} |Y\rangle$, where $|Y\rangle$ is the
GNS vector of the discounted payoff. The real $L^2$ representation makes the
realness and positivity properties manifest.
\end{corollary}

% =======================================================================
% §4.3–4.5 (Semiclassical Expansion, Comparison, Optimal Hedge, Bloch)
% MIGRATED TO ONLINE SUPPLEMENT §S.1 for ≤35 pp main paper target (v4.1)
% =======================================================================

\subsection{The Semiclassical Expansion and the $\mathcal{O}(\tau^{-1})$ Skew}
\label{sec:semiclassical}

A semiclassical (WKB) expansion in $\hbar_{\mathrm{econ}}$ yields tractable
closed-form corrections to the implied volatility surface. The detailed derivation
is given in Online Supplement~\S S.1; we state the main result here.

\begin{assumption}[Semiclassical regime]\label{ass:wkb}
The economic Planck constant satisfies $\hbar_{\mathrm{econ}} \ll 1$, and the
market state $\rho$ is a perturbation of a classical state. Specifically, we assume
that the Wigner function $W_\rho(x, p)$ of the market state admits an expansion
\begin{equation}\label{eq:wigner-expansion}
    W_\rho(x, p) = W_{\mathrm{cl}}(x, p) + \hbar_{\mathrm{econ}} W_1(x, p)
    + \hbar_{\mathrm{econ}}^2 W_2(x, p) + \mathcal{O}(\hbar_{\mathrm{econ}}^3),
\end{equation}
where $W_{\mathrm{cl}}(x, p)$ is a classical probability density on the phase space
$(x, p) \in \mathbb{R}^2$.
\end{assumption}

Under this assumption, the expectation of any observable $A \in \mathcal{A}$ can
be expanded as
\begin{equation*}
    \omega(A) = \int_{\mathbb{R}^2} a_{\mathrm{cl}}(x, p) W_{\mathrm{cl}}(x, p)
    \, dx \, dp + \mathcal{O}(\hbar_{\mathrm{econ}}),
\end{equation*}
where $a_{\mathrm{cl}}(x, p)$ is the classical symbol of $A$. This is the
non-commutative analogue of the method of characteristics for classical stochastic
processes.

\begin{theorem}[Effective implied volatility]\label{thm:effective-vol}
Under Assumption \ref{ass:wkb}, the implied volatility $\sigma_{\mathrm{imp}}$ of
a European option with log-moneyness $x = \log(S_0/K)$ (consistent with §1,
eq.~\eqref{eq:skew-expansion-intro}) and time-to-maturity $\tau$ admits the
short-maturity expansion
\begin{equation}\label{eq:effective-vol}
    \sigma_{\mathrm{imp}}^2(x, \tau) = \sigma_0^2
    + \hbar_{\mathrm{econ}} \cdot \gamma \cdot \frac{x}{\tau}
    + \hbar_{\mathrm{econ}}^2 \cdot \beta \cdot \frac{1}{\tau^2}
    + \mathcal{O}(\hbar_{\mathrm{econ}}^3),
\end{equation}
where $\sigma_0^2 = \omega(\hat{P}^2)$ is the classical variance,
$\gamma = \omega([\hat{X}, \hat{P}]_+)$ is the skew coefficient, and
$\beta = \omega(\hat{X}^2)$ is the convexity coefficient.
\end{theorem}

\begin{remark}[Skew sign convention and comparison with classical models]
\label{rem:skew-sign}
The skew coefficient $\gamma = \omega([\hat{X}, \hat{P}]_+) > 0$ in the equity
leverage-effect regime. The quantum model produces an $\mathcal{O}(\tau^{-1})$
short-maturity skew, faster than BSM (no skew), Heston ($\mathcal{O}(1)$ finite),
and rough volatility ($\mathcal{O}(\tau^{H-1/2})$ for $H < 1/2$). Full proofs,
the model-comparison table, and the Bloch sphere worked example are in Online
Supplement~S.1. Empirical calibration is in Section~\ref{sec:empirical-calibration}.
\end{remark}

\begin{proof}
The proof proceeds via a Wigner--Moyal phase-space expansion in $\hbar_{\mathrm{econ}}$
followed by term-by-term identification of the implied-volatility corrections. The
detailed three-step calculation is given in Online Supplement~S.1.1.
\end{proof}

\begin{corollary}[Skew explosion]\label{cor:skew-explosion}
For any $\hbar_{\mathrm{econ}} > 0$,
\begin{equation*}
    \lim_{\tau \to 0}\,\frac{\partial \sigma^2_{\mathrm{imp}}}{\partial x}
    = \lim_{\tau \to 0}\,\frac{\hbar_{\mathrm{econ}}\gamma}{\tau}
    = +\infty.
\end{equation*}
Consequently, by the chain rule
$\partial\sigma_{\mathrm{imp}}/\partial x
= (2\sigma_{\mathrm{imp}})^{-1}\,\partial\sigma^2_{\mathrm{imp}}/\partial x$,
we also have
$\lim_{\tau\to 0}\partial\sigma_{\mathrm{imp}}/\partial x = +\infty$
whenever $\sigma_{\mathrm{imp}} \to \sigma_0 > 0$.
\end{corollary}

\begin{proposition}[Quantum optimal hedge]\label{prop:quantum-optimal-hedge}
For a European call with log-moneyness $x$ and time-to-maturity $\tau$, the quantum
optimal hedge in the single-qubit parametrization is
\begin{equation}\label{eq:quantum-optimal-hedge}
    \hat{\Pi}^* = \Delta_{\mathrm{BS}}(\sigma_0)
    + \hbar_{\mathrm{econ}} \cdot \Gamma_{\mathrm{BS}}(\sigma_0) \cdot \gamma
      \cdot \frac{x}{\tau}
    + \mathcal{O}(\hbar_{\mathrm{econ}}^2),
\end{equation}
where $\Delta_{\mathrm{BS}}$ is the BSM delta at classical volatility $\sigma_0$
and $\Gamma_{\mathrm{BS}}$ is the BSM gamma.
\end{proposition}
